\documentclass[10pt]{article}
\usepackage[dvips]{graphicx}
%%%%%%%%%%%%%%%%%%%%%%%%%%%%%%%%%%%%%%%%%%%%%%% 
\usepackage{eucal}
\usepackage{amssymb,amsmath}
\usepackage{booktabs}
\usepackage{citesort}
\usepackage{url}
\usepackage{multirow}
%%%%%%%%%%%%%%%%%%%%%%%%%%%%%%%%%%%%%%%%%%%%%%%
\begin{document}

\baselineskip=20pt

\begin{center}
{\bf \huge Using \textsf{Q-SPECTOR}}
  
\vskip 30pt
 
\noindent
by Vladimir Pelmenschikov$^{*}$, David A. Case, and Louis Noodleman 
 
\vskip 10pt
 
\noindent
Contribution from the Department of Molecular Biology, TPC15, The
Scripps Research Institute, 10550 North Torrey Pines Road, La Jolla,
California 92037

\vskip 10pt

\today

\vskip 10pt

\end {center}

\noindent $^{*}$contact author; e-mail: vovan@scripps.edu

\vskip 30pt
 
\noindent
{\bf Abstract:} \textsf{Q-SPECTOR} is a Python program to analyse vibrational outputs from 
several mainstream computational chemistry software. \textsf{Q-SPECTOR} 
can build spectra based on the computed intensities (IR and VCD). \textsf{Q-SPECTOR} also generates 
PVDOS (Partial Vibrational Density od States) and {\it approximate} KED (Kinetic Energy Distribution) spectra based on 
cartesian displacements of the normal modes. The current version \textsf{Q-SPECTOR v1.5} KED function detects bond 
stretching and valence angle bending motions, and torsional angle detection is in the development. 
The program has very flexible interface for PVDOS and KED atoms selection (based on atomic labels, weights, or elements). 
'Stick' spectra are broadened using the pseudo-Voigt profile, 
with limiting cases of Gaussian or Lorentz broadening available as well. The current version has \textsf{GAUSSIAN} and 
\textsf{AMBER} vibrational output reading modules; support for \textsf{ADF} and \textsf{JAGUAR} is under the way. 

\newpage
 
\noindent
{\bf I. Introduction}
..no internal coordinates mess... in MM the force constants are defined, and in QC there are no bond stretch or angle bending 
force constants as such.

\vskip 30pt 

\noindent
{\bf II. Partial Vibrational Density of States (PVDOS) Analysis.}

\vskip 10pt
PVDOS can be characterized via the corresponding mode composition factor $e^2_{j\alpha}$. By the definition, $e^2_{j\alpha}$ is equal to the $<T_{j\alpha}>$ fraction of the ($<...>$ averaged) kinetic energy stored in motion of atom $j$ $(=1,2,...,N)$ in the normal vibrational mode $\alpha$ $(=1,2,...,3N-6)$:
\begin{equation}
\label{eq:modecompfactordef}
e^2_{j\alpha} = \frac{<T_{j\alpha}>}{\sum<T_{i\alpha}>} 
\end{equation}
The mode composition factor is related to the atomic masses $m_j$ and normal mode displacement vectors $\boldsymbol{r}_{j\alpha}$ 
via the well-known formula: 
%%%%%%%%%%%%%%%%%%%%%%%%%%%%%%%%
\begin{equation}
\label{eq:modecompfactor}
e^2_{j\alpha} = \frac{m_j \boldsymbol{r}^2_{j\alpha}}{\sum m_i\boldsymbol{r}^2_{i\alpha}} 
\end{equation}
%%%%%%%%%%%%%%%%%%%%%%%%%%%%%%%%
The above equation \ref{eq:modecompfactor} can be easily proven using the following three relations of the classic harmonic motion.
First, the average kinetic energy $<T_{j\alpha}>$ is equal to the average potential energy $<U_{j\alpha}>$ according to the virial theorem, 
thus being a half of the full energy of the oscillator $E_{j\alpha}$ $(= <E_{j\alpha}>)$ :
%%%%%%%%%%%%%%%%%%%%%%%%%%%%%%%%
\begin{equation}
\label{eq:virial}
<T_{j\alpha}> = <U_{j\alpha}> = \frac{1}{2} E_{j\alpha}
\end{equation}
%%%%%%%%%%%%%%%%%%%%%%%%%%%%%%%%
Second, the full energy of the oscillator $E_{j\alpha}$ is:
%%%%%%%%%%%%%%%%%%%%%%%%%%%%%%%%
\begin{equation}
\label{eq:full}
E_{j\alpha} = U_{j\alpha}\Big\lvert_{T_{j\alpha} = 0} = \frac{1}{2} k_{j\alpha} \boldsymbol{r}^2_{j\alpha}  
\end{equation}
%%%%%%%%%%%%%%%%%%%%%%%%%%%%%%%%
where $k_{j\alpha}$ is the force constant for the motion of atom $j$ in mode $\alpha$. Third, $k_{j\alpha}$ and $m_j$ of atom $j$ are 
connected via the frequency of the mode, $f_\alpha$:
%%%%%%%%%%%%%%%%%%%%%%%%%%%%%%%%
\begin{equation}
\label{eq:freq}
f_\alpha = \frac{1}{2\pi} \sqrt{\frac{k_{j\alpha}}{m_j}}
\end{equation} 
%%%%%%%%%%%%%%%%%%%%%%%%%%%%%%%%
Collecting equations \ref{eq:modecompfactordef} and \ref{eq:virial} through \ref{eq:freq}, we obtain the neccessary relation \ref{eq:modecompfactor}
used to calculate PVDOS in \textsf{Q-SPECTOR}.  

\vskip 30pt 
 
\noindent
{\bf III. Kinetic Energy Distribution (KED) Analysis}

\vskip 5pt 

Notably, atomic PVDOS considered above is essentially equivalent to the single atom KED. However, KED is normally considered as kinetic energy stored in a stretching,
bending, or torsional vibration. Calculating kinetic energies associated with these types of oscillations would involve coupling displacements of at least two atoms. 
Consideration of a covalently bonded atom pair (either as a fragment of a large molecule or in a diatomic molecule) 
will be essential for the stretching and bending KED analysis following in subsections below.   

\vskip 5pt

{\bf III.a. Atom pairs.} For the two atoms 1 and 2, the displacement vector 
$\boldsymbol{r}_{\textrm{CM}}$ of the Center of Mass of the 1-2 pair (CM) is given via formula (here, the normal mode index is omitted for clarity):
%%%%%%%%%%%%%%%%%%%%%%%%%%%%%%%%
\begin{equation}
\label{eq:CM}       
\boldsymbol{r}_{\textrm{1-2}} = \frac{m_1\boldsymbol{r}_{1} + m_2\boldsymbol{r}_{2}}{m_1 + m_2}   
\end{equation} 
%%%%%%%%%%%%%%%%%%%%%%%%%%%%%%%%
Using \ref{eq:CM}, atom 1 and 2 displacement vectors $\boldsymbol{r}_{1_{\textrm{1-2}}}$ and 
$\boldsymbol{r}_{2_{\textrm{1-2}}}$ in the system moving with the center of mass of 1-2 pair are:
%%%%%%%%%%%%%%%%%%%%%%%%%%%%%%%%
\begin{eqnarray}
\label{eq:r12CM}
\boldsymbol{r}_{1_{\textrm{1-2}}} =  \boldsymbol{r}_{1} - \boldsymbol{r}_{\textrm{1-2}} =  
\frac{m_2(\boldsymbol{r}_{1} - \boldsymbol{r}_{2})}{m_1 + m_2} \\
\boldsymbol{r}_{2_{\textrm{1-2}}}  =  \boldsymbol{r}_{2} - \boldsymbol{r}_{\textrm{1-2}} =  
\frac{m_1(\boldsymbol{r}_{2} - \boldsymbol{r}_{1})}{m_1 + m_2}\nonumber 
\end{eqnarray}
%%%%%%%%%%%%%%%%%%%%%%%%%%%%%%%%
Finally, according to equation \ref{eq:r12CM} the collective KED factor associated with the two atoms oscillations relatively to the center of mass 
would be (cf equation \ref{eq:modecompfactor}):
%%%%%%%%%%%%%%%%%%%%%%%%%%%%%%%%
\begin{align}
\label{eq:12KED}
m_1\boldsymbol{r}^2_{1_{\textrm{1-2}}} + m_2\boldsymbol{r}^2_{2_{\textrm{1-2}}} & = 
\frac{m_1 m_2}{m_1 + m_2} \boldsymbol{r}_{12}^2  \\
\textrm{where} \ \boldsymbol{r}_{12} & = \boldsymbol{r}_{1} - \boldsymbol{r}_{2}\nonumber 
\end{align}  
%%%%%%%%%%%%%%%%%%%%%%%%%%%%%%%%
The importance of equation \ref{eq:12KED} will be obvious from the following subsections. 

\vskip 5pt

%%%%%%%%%%%%%%%%%%%%%%%%%%%%%%%%
\begin{figure}[t]
\centering{\includegraphics[width=0.5\textwidth,angle=0]
{/home/vovan/EDIT/Q-SPECTOR/stretch-labels.eps}}
\caption{A schematic representation of $\boldsymbol{r}_{1_{\textrm{1-2}}\parallel}$ and $\boldsymbol{r}_{2_{\textrm{1-2}}\parallel}$ 
displacement vectors
used to calculate a bond stretch KED in \textsf{Q-SPECTOR}. Vectors $\boldsymbol{R}_{12}$ and $\boldsymbol{R}_{1^{\prime}2^{\prime}}$ 
connect the two atoms in their equilibrium positions (points 1 and 2) and after the displacement is applied (points $1^{\prime}$ and $2^{\prime}$).
All the vectors are given relatively to the 1-2 ($1^{\prime}$-$2^{\prime}$) center of mass reference system.}
\label{fig:stretch}
\end{figure}
%%%%%%%%%%%%%%%%%%%%%%%%%%%%%%%%
{\bf III.b. Valence Bond Stretching KED.} 
Here we will do the following approximation for the bond stretching: 
the kinetic energy stored in 1-2 stretch vibration (assuming the two atoms 1 and 2 covalently bound) results from {\it projections} 
$\boldsymbol{r}_{1_{\textrm{1-2}}\parallel}$ and $\boldsymbol{r}_{2_{\textrm{1-2}}\parallel}$ of their
displacement vectors $\boldsymbol{r}_{1_{\textrm{1-2}}}$ and $\boldsymbol{r}_{2_{\textrm{1-2}}}$ on the vector $\boldsymbol{R}_{12}$ 
connecting these two atoms equilibrium positions, see Figure \ref{fig:stretch}. Notably, the displacements are considered relatively to the 1-2 pair center of mass. 
The bond stretch kinetic energy is then proportional to the sum of 1 and 2 kinetic energies (cf equation \ref{eq:modecompfactor}) resulting from 
the $\boldsymbol{r}_{1_{\textrm{1-2}}\parallel}$ and $\boldsymbol{r}_{2_{\textrm{1-2}}\parallel}$ displacements:
%%%%%%%%%%%%%%%%%%%%%%%%%%%%%%%%
\begin{equation}
\label{eq:12stretchKED}
m_1\boldsymbol{r}^2_{1_{\textrm{1-2}}\parallel} + m_2\boldsymbol{r}^2_{2_{\textrm{1-2}}\parallel} = 
\frac{m_1 m_2}{m_1 + m_2}\boldsymbol{r}_{12\parallel}^2
\end{equation}
%%%%%%%%%%%%%%%%%%%%%%%%%%%%%%%%
where $\boldsymbol{r}_{12\parallel}$ is the projection of $\boldsymbol{r}_{1}-\boldsymbol{r}_{2}$ 
normal mode displacement vectors difference onto $\boldsymbol{R}_{12}$. 
The right side of equation \ref{eq:12stretchKED} is derived using equations \ref{eq:r12CM} and \ref{eq:12KED}. Finally, 
using equations \ref{eq:modecompfactor} and \ref{eq:12stretchKED}, 
1-2 stretch KED composition factor $e^2_{12\alpha}$ for the mode index $\alpha$ is (the notation is analogous to $e^2_{j\alpha}$):
%%%%%%%%%%%%%%%%%%%%%%%%%%%%%%%%
\begin{equation}
\label{eq:KEDstretchfactor}
e^2_{12\alpha} = \frac{\frac{m_1 m_2}{m_1 + m_2}\boldsymbol{r}_{12\alpha\parallel}^2}{\sum m_i\boldsymbol{r}^2_{i\alpha}} 
\end{equation}   
%%%%%%%%%%%%%%%%%%%%%%%%%%%%%%%%
Equation \ref{eq:KEDstretchfactor} is used to calculate 1-2 valence bond stretch KED in \textsf{Q-SPECTOR}. The list of 
$e^2_{12\alpha}$ values obtained for all $\alpha$ can be used to analyse where 1-2 stretch vibrations occur along the vibraional 
modes range, and what is their relative weight. To generate KED stretch spectra, \textsf{Q-SPECTOR} applies broadening to $e^2_{12\alpha}$ list.
It is easy to show that equation \ref{eq:KEDstretchfactor} results in $e^2_{12}$ = 1 for a diatomic molecule vibrational mode, 
using $\boldsymbol{r}_{12\parallel} =\boldsymbol{r}_{12}$ and $\boldsymbol{r}_2 = -\frac{m_1}{m_2} \boldsymbol{r}_1$. 
This reflects proper behavior of formula \ref{eq:KEDstretchfactor} in case when 100$\%$ of the mode vibrational energy is stored in 1-2 stretch. 

\vskip 5pt

%%%%%%%%%%%%%%%%%%%%%%%%%%%%%%%%
\begin{figure}[t]
\centering{\includegraphics[width=\textwidth,angle=0]
{/home/vovan/EDIT/Q-SPECTOR/bend.eps}}
\caption{{\bf a}: a schematic representation of $\boldsymbol{r}_{1_{\textrm{1-2}}\parallel}$ and $\boldsymbol{r}_{2_{\textrm{1-2}}\parallel}$ 
displacement vectors used to calculate valence and torsional angle bendings in \textsf{Q-SPECTOR}. {\bf b}: a schematic representation of 
$\boldsymbol{r}_{1_{\textrm{1-2}}\perp_{{\textrm{1-2-3}}\parallel}}$ displacement contribution to 1-2-3 angle bending.}
\label{fig:bend}
\end{figure}
%%%%%%%%%%%%%%%%%%%%%%%%%%%%%%%%
{\bf III.b. Valence Angle Bending KED.} 
Describing 1-2-3 valence angle bending, where 1-2 and 2-3 atom pairs are assumed to be bound, 
we first naturally suppose that bending (in contrast to stretching, see above) comes from the displacements, perpendicular to the 1-2 and 2-3 bonding vectors 
$\boldsymbol{R}_{21}$ and $\boldsymbol{R}_{23}$. These normal displacements will be denoted $\boldsymbol{r}_{1_{\textrm{1-2}}\perp}$ and $\boldsymbol{r}_{2_{\textrm{1-2}}\perp}$ for 1-2 pair, see Figure \ref{fig:bend}, 
and $\boldsymbol{r}_{2_{\textrm{2-3}}\perp}$ and $\boldsymbol{r}_{3_{\textrm{2-3}}\perp}$ for 2-3 pair. Notably, 
%%%%%%%%%%%%%%%%%%%%%%%%%%%%%%%%
\begin{equation}
\label{eq:parallel+perp}
\boldsymbol{r}_{i_{i-j}} = \boldsymbol{r}_{i_{i-j}\parallel} + \boldsymbol{r}_{i_{i-j}\perp}
\end{equation}   
%%%%%%%%%%%%%%%%%%%%%%%%%%%%%%%%
1-2-3 valence angle bending is also limited to the displacements which occur within the plane formed by the three atoms 1-2-3. One more 
The second step of projection should be done, removing those components of displacements which are normal to the 1-2-3 plane 
(the normal components can be used for the description of torsional contributions to KED, see below). 
Following our notation scheme, $\boldsymbol{r}_{1_{\textrm{1-2}}\perp}$
when projected onto 1-2-3 plane would be called $\boldsymbol{r}_{1_{\textrm{1-2}}\perp_{{\textrm{1-2-3}}\parallel}}$, see Figure \ref{fig:bend}.  
Then the amount of kinetic energy stored in 1-2-3 angle bending motion can be approximated as sum of 1-2 and 2-3 'bending' 
components (cf. equation \ref{eq:12stretchKED}), given below:
\begin{eqnarray}
\label{eq:bendKED}
m_1\boldsymbol{r}^2_{1_{\textrm{1-2}}\perp_{{\textrm{1-2-3}}\parallel}} + m_2\boldsymbol{r}^2_{2_{\textrm{1-2}}\perp_{{\textrm{1-2-3}}\parallel}} = 
\frac{m_1 m_2}{m_1 + m_2} \boldsymbol{r}_{12\perp_{\textrm{1-2-3}\parallel}}^2\\
m_2\boldsymbol{r}^2_{2_{\textrm{2-3}}\perp_{{\textrm{1-2-3}}\parallel}} + m_3\boldsymbol{r}^2_{3_{\textrm{2-3}}\perp_{{\textrm{1-2-3}}\parallel}} = 
\frac{m_2 m_3}{m_2 + m_3} \boldsymbol{r}_{23\perp_{\textrm{1-2-3}\parallel}}^2\nonumber
\end{eqnarray}
where $\boldsymbol{r}_{12\perp_{\textrm{1-2-3}\parallel}}$ is the component of $\boldsymbol{r}_{1}-\boldsymbol{r}_{2}$
normal mode displacement vectors difference normal to $\boldsymbol{R}_{12}$, further projected onto the 1-2-3 plane.   
For mode index $\alpha$, the fraction of kinetic energy stored in 1-2-3 angle bending then can be written as follows :
%%%%%%%%%%%%%%%%%%%%%%%%%%%%%%%%
\begin{equation}
\label{eq:KEDbendfactor}
e^2_{123\alpha} = \frac{\frac{m_1 m_2}{m_1 + m_2}\boldsymbol{r}_{12\alpha\perp_{\textrm{1-2-3}\parallel}}^2 + 
\frac{m_2 m_3}{m_2 + m_3}\boldsymbol{r}_{23\alpha\perp_{\textrm{1-2-3}\parallel}}^2}{\sum m_i\boldsymbol{r}^2_{i\alpha}} 
\end{equation}   
%%%%%%%%%%%%%%%%%%%%%%%%%%%%%%%%
...
\vskip 30pt 

\newpage 

\noindent
{\bf IV. \textsf{Q-SPECTOR} Input File and Usage.} The input (instructions) file format is as follows:

\vskip 5pt 

\texttt{
\begin{tabular}{l}
\&output\\
spectra [IR,VCD,PVDOS,KEDSTRETCH,KEDBEND]\\
shape [SIGMA GAMMA]\\ 
range START STOP\\
resolution RES\\
\&\\
\&selection\\
element1 [E1\_1 E1\_2 ...]\\
element2 [E2\_1 E2\_2 ...]\\
element3 [E3\_1 E3\_2 ...]\\
weight1 [W1\_1 W1\_2 ...]\\
weight2 [W2\_1 W2\_2 ...]\\
weight3 [W3\_1 W3\_2 ...]\\
label1 [L1\_1 L1\_2 ...]\\
label2 [L2\_1 L2\_2 ...]\\
label3 [L3\_1 L3\_2 ...]\\
cutoff CUTOFF\\
\&\\
\end{tabular}
}
\vskip 20pt
\noindent
...the above format is to be explained...
 
\vskip 5pt

\noindent
Using \textsf{Q-SPECTOR} on input.qsp with \textsf{GAUSSIAN} outputfile :\\
\texttt{>Q-Spector-1.5.4.py GAUSSIAN input.qsp output.g03out output.dat}\\ 
which results in output.dat spectrum data file. The output data has two column format, where the first column is the wavenumber (cm$^{-1}$) and the second is intensity.  

\vskip 30pt

\noindent
{\bf Appendix}

\vskip 5pt 

Exact formulae for projected displacement differences $\boldsymbol{r}_{12\parallel}$ and $\boldsymbol{r}_{12\perp_{\textrm{1-2-3}\parallel}}$ 
used to calculate $e^2_{12\alpha}$ (equation \ref{eq:KEDstretchfactor}) and $e^2_{123\alpha}$ (equation \ref{eq:KEDbendfactor}), correspondingly.
...


%%%%%%%%%%%%%%%%%%%%%%%%%%%%%%%%%%%%%%%%%
%\bibliographystyle{achemso}
\bibliographystyle{jpc}
\bibliography{nitrogenase}
%%%%%%%%%%%%%%%%%%%%%%%%%%%%%%%%%%%%%%%%%

\end{document}